% PaperLab venue template: JOSS (Journal of Open Source Software)
% Minimal compilable skeleton — article class, common packages only.
% Note: real JOSS uses the `dgemc`/`agujournaltemplate` class at submission time;
% this skeleton keeps the paper compilable everywhere with zero exotic packages.
\documentclass[11pt]{article}
\usepackage[a4paper,margin=1in]{geometry}
\usepackage{graphicx}
\usepackage{booktabs}
\usepackage{natbib}
\usepackage{hyperref}

\title{TITLE: software-name — one-line statement of what it does}
\author{AUTHOR NAME\thanks{Affiliation; email@example.com} \and AUTHOR TWO\thanks{Affiliation}}
\date{\today}

\begin{document}
\maketitle

\begin{abstract}
\noindent PLACEHOLDER-ABSTRACT-SUMMARY: JOSS papers are short (500--1000 words) descriptions of
research software with existing academic use. State the problem, the software,
and the referenceable DOI release. Reviewers check installation, usage, and tests,
not novelty.
\\[2mm]
\textbf{Keywords:} keyword one; keyword two
\end{abstract}

\section{Summary}
PLACEHOLDER-SUMMARY: the problem domain, why existing tools fall short, and what
this software provides.

\section{Statement of need}
PLACEHOLDER-NEED: who uses it, for what task, and the gap it fills
(see Figure~\ref{fig:placeholder}).

\begin{figure}[t]
\centering
% \includegraphics[width=0.85\linewidth]{figures/figure1.pdf}
\caption{PLACEHOLDER-CAPTION: overview of the software and its workflow.}
\label{fig:placeholder}
\end{figure}

\section{Availability}
PLACEHOLDER-AVAILABILITY: repository URL, archived DOI (Zenodo), license,
and the exact version reviewed. State the test suite and CI status.

\section{References}
\bibliographystyle{plainnat}
\bibliography{refs}

\end{document}
